\printedpage{39}
% head: A. Grothendieck : Dérivateurs, Ch. I, Édition des Universités de Montpellier II et Paris VII
% foot: Version 8 juillet 2001
\msp{page 67}

\noindent Thus the arrow under consideration $(**)$ is interpreted as
\[
H_*^{\D}(y\backslash X, M_X) \longrightarrow M
\]
and this is the adjunction homomorphism $p_!(p^*(M)) \to M$.
\[
\begin{tikzcd}
X \arrow[rr, "f"] \arrow[dr, "p"{swap}] & & Y \arrow[dl, "q"] \\
& e &
\end{tikzcd}
\]

\noindent Thus it is an isomorphism for every $M$ if and only if $y\backslash X$ is $\D$-contractible.

\noindent We are going to state what we have proved.

\noindent\textbf{Theorem 1.} \textit{Let}
\[
\begin{tikzcd}
X \arrow[rr, "f"] \arrow[dr, "p"{swap}] & & Y \arrow[dl, "q"] \\
& e &
\end{tikzcd}
\]
\textit{be in} $\Dia$. \textit{We assume that} $\D$ \textit{is a prederivator satisfying} \textbf{Der 2}, \textit{that} $p_*, q_*$ \textit{and} $f_!$ \textit{exist, that} $f_!$ \textit{satisfies} \textbf{Der 4} (d) \textit{for the calculation of fibers. Equivalent conditions:}

a) $f$ \textit{is a costrict} $\D$\textit{-equivalence, i.e.} (prop. 2) \textit{for every} $y \in \Ob\, Y$, $y\backslash X$ \textit{is} $\D$\textit{-contractible;}

b) \textit{for every} $F \in \Ob\, \D(Y)$, \textit{the following homomorphism is} \texttt{[an]} \textit{isomorphism}
\[
H^*_{\D}(Y, F) \isoto H^*_{\D}(X, f^*(F)) \ ;
\]

c) \textit{for every} $M \in \Ob(\mathcal{A} = \D(e))$, \textit{the adjunction homomorphism is} \texttt{[an]} \textit{isomorphism}
\[
f_!(M_X) = f_!(f^*(M_Y)) \isoto M_Y \ .
\]

I am going to renumber the axioms on the prederivators. % [sic] le prédérivateurs -> les prédérivateurs, article singulier devant un pluriel

\noindent\textbf{Der 1}, \textbf{Der 2} as before (p. 45).

\noindent\textbf{Der 3}\footnote[28]{Axiom of aspheric indexing categories, \textit{cf.} variants \textbf{Der 3} a) b) c) page 97, and especially pages 111-112.} : Let $X$ be in $\Dia$, $p : X \to e$. Suppose $X$ contractible.\footnote[29]{N.B. It suffices to assume that $X$ has a final object. But this is not a self-dual axiom.} Then
\[
p^* : M \longmapsto M_X \quad , \qquad p^* : \mathcal{A} = \D(e) \longrightarrow \D(X)
\]
is fully faithful.

\printedpage{40}
% head: A. Grothendieck : Dérivateurs, Ch. I, Édition des Universités de Montpellier II et Paris VII
% foot: Version 8 juillet 2001
\msp{page 68}

\noindent\textbf{Der 4} (g) : existence of the $f_*$.

\noindent\textbf{Der 4} (d) : existence of the $f_!$.

\noindent\textbf{Der 5} (g) Calculation of the fibers of $f_*(F)$.

\noindent\textbf{Der 5} (d) Calculation of the fibers of $f_!(F)$.

\noindent The old axiom \textbf{Der 5} (g) follows from \textbf{Der 2}, \textbf{Der 3}, \textbf{Der 4} (g,d), \textbf{Der 5} (d), thanks to the preceding theorem.

An interesting aspect of the preceding theorem, if we view it as stating the equivalence of a) and b), is that this statement does not appeal to homological notions, but that its proof does appeal to them, in a relatively strong form (calculation of the fibers of $f_!(F)$). For the base change theorems for relative $\D$-cohomology, \textit{i.e.} for the $f_*^{\D}$, we also constantly need the calculation of the fibers of $f_*(F)$. Thus it seems that we cannot dispense with any of the preceding axioms, including \textbf{Der 4} (d) and \textbf{Der 5} (d), even if we only want to concern ourselves with the formalism of the $f^*, f_*$ (cohomological aspect). In what follows, we shall therefore assume these axioms satisfied.

\noindent\textbf{Theorem 2.} \textit{Let} $f : X \to Y$ \textit{be in} $\Dia$. \textit{Equivalent conditions:}

a) \textit{For every} $y_0 \in \Ob\, Y$, \textit{the canonical functor}
\[
X_{y_0} \overset{\varphi}{\longrightarrow} X/y_0 \quad (\simeq X \times_Y Y/y_0) \ ,
\]
\textit{where} $X_{y_0}$ \textit{denotes the fiber at} $y_0$, \textit{is a co-strict} $\D$\textit{-equivalence.}

b) \textit{For every} $x_0 \in \Ob\, X$, \textit{the induced functor on the categories colocalized at} $x_0$ \textit{and at} $y_0 = f(x_0)$
\[
x_0\backslash X \longrightarrow y_0\backslash Y
\]
\textit{has} $\D$\textit{-contractible fibers.}

\msp{page 69}

c) \textit{The formation of the} $f_*F$ \textit{(}$F$ \textit{in} $\D(X)$\textit{) commutes with every base change} $Y' \to Y$, \textit{and this remains true after every base change -- i.e. if we have} $Y'' \to Y' \to Y$, \textit{whence}
\[
\begin{tikzcd}
X'' \arrow[r] \arrow[d, "f''"{swap}] & X' \arrow[r] \arrow[d, "f'"] & X \arrow[d, "f"] \\
Y'' \arrow[r, "p'"{swap}] & Y' \arrow[r, "p"{swap}] & Y
\end{tikzcd}
\]
\textit{the ``base change theorem'' is true for} $f'_*$ \textit{and} $p'^{*}$.

c$'$) \textit{The formation of the} $f_*(F)$ \textit{(}$F \in \Ob\, \D(X)$\textit{) commutes with passage to the fibers and this remains true after every base change} $p : Y' \to Y$.

c$''$) \textit{As in} c$'$) \textit{with} $Y' \to Y$ \textit{a localization morphism} $Y/y \to Y$.

d) \textit{If} $Y'' \overset{p}{\to} Y' \to Y$, \textit{with} $p$ \textit{coaspheric, then} $q : X'' = X \times_Y Y'' \to X' = X \times_Y Y'$ \textit{is coaspheric.}

\printedpage{41}
% head: A. Grothendieck : Dérivateurs, Ch. I, Édition des Universités de Montpellier II et Paris VII
% foot: Version 8 juillet 2001

d$'$) \textit{As in} d) \textit{with} $Y'' = e$ (so $Y'$ is a category with final object $e$).

d$''$) \textit{As in} d) \textit{with} $Y'' = \{1\}$, $Y' = \boldsymbol{\Delta}_1$.

.

\noindent\textsc{Proof}. a) $\Leftrightarrow$ b) cf. proposition page 53. Let us show that these conditions imply c). Since they are stable under base change (\textit{cf.} immediate proof page 55) therefore by \textbf{Der 2} we are reduced to proving that $H^*_{\D}(X/y, F) \overset{\sim}{\to} H^*_{\D}(Y_y, F)$ $(F \in \Ob\, \D(Y))$, which follows from theorem 1.\footnote[30]{Thus we use \textbf{Der 5} (d).} (Cf. proof of the theorem page 56, which works with more general hypotheses.)

We have c) $\Leftrightarrow$ c$'$) $\Rightarrow$ c$''$) (the equivalence c) $\Leftrightarrow$ c$'$) by \textbf{Der 2}).

Let us prove c$''$) $\Rightarrow$ a).

\msp{page 70}

\noindent Consider the diagram with cartesian squares
\[
\begin{tikzcd}
X_{y_0} \arrow[r, "q"] \arrow[d, "p"{swap}] & X/y_0 \arrow[r] \arrow[d, "f_0"] & X \arrow[d, "f"] \\
\{y_0\} \arrow[r] & Y/y_0 \arrow[r] & Y
\end{tikzcd}
\qquad\qquad .
\]

\noindent Hypothesis c$''$) proves that for every $F$ in $\D(X/y_0)$, we have $H^*_{\D}(X/y_0, F) \overset{\sim}{\to} H^*_{\D}(X_{y_0}, F)$. By theorem 1, this implies that $f$ is a co-strict $\D$-equivalence. % [sic] un D-équivalence -> une D-équivalence, accord du genre

\noindent\textbf{Corollary 1.} Dual statement of the theorem, for the functors $X_{y_0} \to y_0\backslash X$, $X/x_0 \to Y/y_0$, and base change for the $j_!(F)$ \ldots

\noindent\textbf{Definition 4.} Under the equivalent conditions of theorem 1, we say that $f$ is $\D$-\textit{proper.} The dual condition is called $f$ $\D$-\textit{smooth.}

\noindent\textbf{Corollary 2.} \textit{Let} $f : X \to Y$ \textit{be in} $\Dia$. \textit{If} $f$ \textit{is} $\D$\textit{-proper or} $\D$\textit{-smooth, and has} $\D$\textit{-contractible fibers, then} $f$ \textit{is a strict and costrict} $\D$\textit{-equivalence. Better,} $f$ \textit{is a} universal $\D$\textit{-equivalence} \textit{i.e. is a} $\D$\textit{-equivalence and remains so after every base change.}\footnote[31]{Generalizes: \textbf{Corollary 2$'$.} \textit{Let}
\[
\begin{tikzcd}[ampersand replacement=\&]
X \arrow[rr, "f"] \arrow[dr] \& \& Y \arrow[dl] \\
\& S \&
\end{tikzcd}
\]
\textit{with} $X, Y$, $\D$\textit{-proper over} $S$. \textit{If} $f$ \textit{induces a} $\D$\textit{-equivalence on the fibers it is a} $\D$\textit{-equivalence. Likewise in the smooth case.}\par
This is proposition 6 below.}

\printedpage{42}
% head: A. Grothendieck : Dérivateurs, Ch. I, Édition des Universités de Montpellier II et Paris VII
% foot: Version 8 juillet 2001
\msp{page 71}

\noindent It suffices to prove that $f$ is a $\D$-equivalence. Suppose for example $f$ $\D$-proper. It suffices to prove that if $Y$ is $\D$-contractible, then $X$ is so (if need be, applying it to the $Y/y_0$ after base change, to show that the $X/y_0$ are contractible). For this, it suffices to prove that $f$ is a $\D$-equivalence, and we know that this is the case if for every $F$ \textit{constant} on $Y$, the homomorphism $F \overset{\text{adj}}{\longrightarrow} f_*f^*F$ is \texttt{[an]} isomorphism. In fact, this is true for \textit{every} $F$.

\noindent\textbf{Corollary 3.} \textit{If} $f : X \to Y$ \textit{is} $\D$\textit{-proper or} $\D$\textit{-smooth, with} $\D$\textit{-contractible fibers, then} $f^* : \D(Y) \to \D(X)$ \textit{is fully faithful, i.e. the adjunction homomorphisms}
\[
F \longrightarrow f_*f^*F \quad , \qquad f_!f^*F \longrightarrow F
\]
\textit{in} $\D(Y)$ \textit{are isomorphisms.}

\noindent It suffices to see this for $f$ proper, and this follows from the calculation of the fibers of $f_*f^*F$ \ldots

\noindent\textbf{Proposition 5.} \textit{The notion of a} $\D$\textit{-proper functor} (resp. $\D$\textit{-smooth}) \textit{is stable under base change,} (trivial), \textit{and under composition} (use \texttt{[the]} corollary 2). \textit{Every closed immersion is} $\D$\textit{-proper, every open immersion} $\D$\textit{-smooth. More generally, every local isomorphism is} $\D$\textit{-smooth, every local co-isomorphism is} $\D$\textit{-proper.}

\noindent\textbf{Proposition 6.} \textit{Let} $X, Y$ \textit{be over} $S$, $f : X \to Y$ \textit{an} $S$\textit{-morphism. Suppose} $X, Y$ \textit{proper over} $S$, \textit{or} $X$ \textit{and} $Y$ \textit{smooth over} $S$, \textit{and that for every} $s \in \Ob\, S$, $f_s : X_s \to Y_s$ \textit{is a} $\D$\textit{-equivalence, then} $f$ \textit{also.}\footnote[32]{\textsc{Proof}. We shall even prove that on the subcategory of coefficients on $Y$ which come from $\D(Y)$ \texttt{[rather from }$\D(S)$\texttt{]}, $f^* : \D(Y) \to \D(X)$ induces a fully faithful functor. If $f$ \texttt{[is]} proper, I interpret it as $q_*q^*F \overset{\sim}{\to} p_*p^*F$ for every $F$, and this follows from the calculation of $p_*$, $q_*$ by passage to the fibers (proper), and from the hypothesis on the $f_s$. OK.}

\msp{page 72}

\noindent\textbf{Theorem 3.} \textit{Let}
\[
\begin{tikzcd}
X \arrow[d, "f"{swap}] & X' \arrow[l, "q"{swap}] \arrow[d, "f'"] \\
Y & Y' \arrow[l, "p"{swap}]
\end{tikzcd}
\]
\textit{be cartesian in} $\Dia$. \textit{Then the base change morphism}
\[
p^*f_* \longrightarrow f'_*q^*
\]
\textit{is an isomorphism if} $f$ \textit{is} $\D$\textit{-proper, or} $p$ $\D$\textit{-smooth.}

\noindent The $\D$-proper case already proved. The $\D$-smooth case: the proof p. 56 is OK.

\printedpage{43}
% head: A. Grothendieck : Dérivateurs, Ch. I, Édition des Universités de Montpellier II et Paris VII
% foot: Version 8 juillet 2001

\subsection*{Comparison of $X$ and $X^{\circ}$.}

Let $X$ be \texttt{[a]} category, we define a category $\Phi(X)$ and functors
\[
\begin{tikzcd}[column sep=small]
& & \Phi(X) \arrow[ddl, "p_X"{swap, name=P}] \arrow[ddr, "q_X"{name=Q}] & & \\
|[alias=L]| \text{\footnotesize\begin{tabular}{c}$\Cat$-cofibrant\\(hence proper)\end{tabular}} & & & & |[alias=R]| \text{\footnotesize\begin{tabular}{c}$\Cat$-fibrant\\(hence smooth)\\ \texttt{[rather }$\Cat$\texttt{-cofibrant}\\ \texttt{(hence proper)]}\end{tabular}} \\
& X & & X^{\circ} &
\arrow[from=P, to=L]
\arrow[from=Q, to=R]
\end{tikzcd}
\]
which will be $\hot$-equivalences, and $\D$-equivalences if $\D$ is a derivator whose domain contains $X$ and $\Phi(X)$. We set
\[
\Ob\, \Phi(X) = \Fl(X) \quad ,
\]
so the same objects as $\mathcal{F}\ell(X) = \Hom(\boldsymbol{\Delta}_1, X)$, but the morphisms are different:

\msp{page 73}
\[
\begin{tikzcd}
Y_1 \arrow[r, "u"] & Y_0 \arrow[d, "v_0"] \\
Y'_1 \arrow[u, "v_1"] \arrow[r, "u'"{swap}] & Y'_0
\end{tikzcd}
\]
a morphism from $u$ to $u'$ is a pair $v_0 : Y_0 \to Y'_0$, $v_1 : Y'_1 \to Y_1$ such that $v_0uv_1 = u'$. Thus source and target $s(u)$ and $b(u)$ are the one contravariant, the other covariant in $u$. We take for $q$ the source functor, for $p$ the target functor. % [sic] l'autre covariant -> l'autre covariante, accord avec « l'une contravariante »

If we apply this to $X^{\circ}$, we find
\[
\begin{tikzcd}
& \Phi(X^{\circ}) \arrow[dl, "p_{X^{\circ}}"{swap}] \arrow[dr, "q_{X^{\circ}}"] & \\
X^{\circ} & & (X^{\circ})^{\circ} = X
\end{tikzcd}
\]
We have $\Ob\, \Phi(X^{\circ}) = \Fl(X^{\circ}) \simeq \Fl(X)$
\[
\begin{tikzcd}
Y_1 \arrow[r, "u"] & Y_0 \arrow[d, "v_1"] & \text{in}\ \ X \\
Y'_1 \arrow[u, "v_0"] \arrow[r, "u'"{swap}] & Y'_0 & \text{in}\ \ X
\end{tikzcd}
\]
and the same homomorphisms in the sense of $\Phi(X)$ and in the sense of $\Phi(X^{\circ})$. Thus
\[
\left\{
\begin{aligned}
&\Phi(X^{\circ}) = \Phi(X) \\
&p_{X^{\circ}} = q_X \\
&q_{X^{\circ}} = p_X \quad ,
\end{aligned}
\right.
\]